\section{Conclusiones}

El presente trabajo ha presentado el diseño, implementación y validación de LPR-Sentinel, un sistema completo de reconocimiento automático de matrículas vehiculares mexicanas basado exclusivamente en datos sintéticos y desplegable sobre hardware de bajo costo. A lo largo del desarrollo se han abordado los principales desafíos que limitan la adopción de esta tecnología en el contexto mexicano: la heterogeneidad de diseños por entidad federativa, la ausencia de bases de datos públicas etiquetadas y las restricciones legales asociadas al uso de imágenes reales. La solución propuesta ha demostrado que es posible construir un pipeline funcional sin necesidad de recopilar ni almacenar datos reales de vehículos o sus propietarios, recurriendo en su lugar a la generación procedimental de matrículas conforme a la Norma Oficial Mexicana NOM-001-SCT-2-2016 y a su posterior aumentación mediante transformaciones que simulan condiciones reales de captura.

El sistema se articula en cinco módulos independientes pero coordinados. El MLP-Generator produce matrículas sintéticas con una fidelidad normativa del 100%, respetando los rangos alfanuméricos asignados a cada estado, las combinaciones cromáticas oficiales y la fuente tipográfica diseñada manualmente a partir de las cotas especificadas en la norma. El MLP-Augmentator aplica 29 transformaciones geométricas, degradaciones ópticas y ajustes fotométricos que permiten generar miles de variantes a partir de una misma matrícula base, introduciendo la diversidad necesaria para entrenar modelos robustos a condiciones adversas. El MLP-Detector, basado en YOLOv11 y exportado a ONNX, localiza la región de la placa en imágenes de escena completa con una precisión media del 94.3% y una latencia de 78 ms en Raspberry Pi 4B. El MLP-Recognizer, basado en una arquitectura Compact Convolutional Transformer y también exportado a ONNX, transcribe el texto de la matrícula con una precisión por carácter del 95.1% en condiciones diurnas y del 91.4% en condiciones nocturnas, con una latencia de 42 ms. Finalmente, el MLP-Pipeline integra todos los componentes anteriores junto con una base de datos SQLite, alcanzando una tasa de éxito integral del 85% en condiciones diurnas y del 70% en condiciones nocturnas sobre un conjunto de prueba real, con una latencia extremo a extremo de 121 ms por imagen.

Las contribuciones más relevantes de este trabajo son, en primer lugar, la demostración de que los datos sintéticos procedimentales, cuando son generados con estricto apego a la normativa y aumentados adecuadamente, pueden sustituir a los datos reales en el entrenamiento de sistemas LPR sin incurrir en problemas éticos o legales. En segundo lugar, la adopción del formato ONNX como interfaz unificada de despliegue permite ejecutar el pipeline completo en dispositivos de bajo costo y consumo energético reducido, como la Raspberry Pi 4B, abriendo la posibilidad de aplicaciones distribuidas y sistemas embebidos. En tercer lugar, la arquitectura modular facilita la mejora independiente de cada componente, la sustitución de modelos subyacentes y la adaptación a futuras actualizaciones de la normativa mexicana sin necesidad de rediseñar el sistema completo. En cuarto lugar, el repositorio público y la documentación proporcionada garantizan la reproducibilidad de los experimentos y permiten a otros investigadores construir sobre esta base.

Las limitaciones del sistema actual incluyen la sensibilidad a rotaciones extremas de la placa (superiores a 60°), a desenfoques de movimiento severos y a condiciones de iluminación muy desfavorables (sombras profundas o contraluz intensa), escenarios en los que la tasa de éxito integral desciende por debajo del 50%. Asimismo, la fuente tipográfica utilizada, aunque respetuosa de las cotas oficiales, no es la fuente exacta empleada por los fabricantes de matrículas mexicanas, lo que podría introducir pequeñas discrepancias en la fase de reconocimiento sobre imágenes reales de muy alta resolución. Finalmente, el conjunto de imágenes reales empleado en la validación, si bien diverso en condiciones de iluminación y ángulos, es limitado en tamaño (200 imágenes) y no cubre la totalidad de los 32 estados ni todos los tipos de vehículos (gubernamentales, diplomáticos, de carga, de personas con discapacidad, etc.).

Las líneas de trabajo futuro se orientan en varias direcciones. En primer lugar, se planea ampliar el conjunto de validación real con imágenes capturadas en los 32 estados mexicanos, recolectadas en colaboración con cuerpos de seguridad y autoridades de tránsito bajo los protocolos éticos correspondientes. En segundo lugar, se explorará la incorporación de técnicas de aprendizaje por transferencia que permitan ajustar los modelos sintéticos con un pequeño número de imágenes reales etiquetadas, mejorando así la generalización a condiciones extremas. En tercer lugar, se desarrollará una versión del pipeline optimizada para aceleración por hardware, empleando la unidad de procesamiento neuronal (NPU) de la Raspberry Pi 5 o de dispositivos como Google Coral TPU, con el objetivo de reducir la latencia por debajo de los 50 ms y alcanzar tasas de 20 fotogramas por segundo. En cuarto lugar, se implementará un módulo de seguimiento temporal (\textit{tracking}) que permita asociar detecciones de la misma placa en secuencias de video, mejorando la fiabilidad del reconocimiento mediante la agregación de múltiples predicciones. Finalmente, se estudiará la extensión del sistema a otros países latinoamericanos con normativas de matriculación similares a la mexicana, adaptando los generadores y configuraciones correspondientes.

En resumen, LPR-Sentinel constituye una prueba de concepto sólida y funcional de que los sistemas de reconocimiento de matrículas pueden ser desarrollados, entrenados y desplegados íntegramente con datos sintéticos, respetando así la privacidad de los ciudadanos y evitando las complejidades legales del uso de imágenes reales. La combinación de generación procedimental, aumentación realista, modelos profundos eficientes y despliegue en formato ONNX sobre hardware de bajo costo ofrece un camino viable para la democratización de esta tecnología, poniéndola al alcance de pequeños municipios, empresas de transporte y comunidades que no disponen de recursos para adquirir soluciones comerciales propietarias. El código fuente, los modelos entrenados y las bases de datos sintéticas generadas se encuentran disponibles en el repositorio público asociado a este trabajo, invitando a la comunidad investigadora a utilizarlos, extenderlos y mejorarlos.

\begin{figure}[h]
    \centering
    \includegraphics[width=0.6\textwidth]{figures/intro.pdf}
    \captionsetup{justification=centering}
    \caption{Example figure. Make sure to generate figures as .pdf files for the best quality.}
    \label{fig:intro}
\end{figure}